\documentclass[10pt,leqno]{amsart}
\usepackage[T1]{fontenc}
\usepackage{lmodern}
\usepackage{amsmath,amssymb}
\usepackage{microtype}
\usepackage[letterpaper,margin=75pt,includeheadfoot]{geometry}
\usepackage[hidelinks]{hyperref}

\newtheorem{theorem}{Theorem}[section]
\newtheorem{corollary}[theorem]{Corollary}
\numberwithin{equation}{section}
\newcommand{\cf}{\operatorname{cf}}
\newcommand{\cl}{\operatorname{cl}}
\let\Small\normalsize
\let\footnotesize\normalsize

\title[A set mapping with small intersections]{A set mapping with small intersections\protect\\ and no large free set}
\author[W. Howe, G. Huffman, J. Kohlmeyer, and L. Kruer]{Will Howe, Grant Huffman, Jensen Kohlmeyer, and Liam Kruer}
\thanks{\raggedright Email addresses: \href{mailto:howew@purdue.edu}{\texttt{howew@purdue.edu}} (W.~Howe);
\href{mailto:huffma97@purdue.edu}{\texttt{huffma97@purdue.edu}} (G.~Huffman);
\href{mailto:jenwin@purdue.edu}{\texttt{jenwin@purdue.edu}} (J.~Kohlmeyer);
\href{mailto:kruerl@purdue.edu}{\texttt{kruerl@purdue.edu}} (L.~Kruer).}
\date{November 10, 2026}
\hypersetup{
  pdftitle={A Set Mapping with Small Intersections and No Large Free Set},
  pdfauthor={Will Howe; Grant Huffman; Jensen Kohlmeyer; Liam Kruer},
  pdfsubject={A set-mapping counterexample at a successor of a singular cardinal},
  pdfcreator={Academic manuscript}
}

\begin{document}
\begin{abstract}
Let $\mu$ be an uncountable cardinal of countable cofinality and suppose that
$2^\mu=\mu^+$. We construct a self-avoiding set mapping on $\mu^+$ whose values
have cardinality at most $\mu$, whose distinct values have intersection of
cardinality less than $\mu$, and which has no free set of cardinality $\mu^+$.
The construction combines a countable neighborhood base with an enumeration
of the $\mu$-sized subsets and a schedule having fewer than $\mu$ requirements
at each bounded depth. In particular, under the generalized continuum
hypothesis, it gives a negative answer to the free-set question of Erd\H{o}s
and Hajnal at $\aleph_{\omega+1}$.
\end{abstract}
\maketitle

\section{Introduction and main result}

Let $S$ be a set. A mapping $F:S\to\mathcal P(S)$ is \emph{self-avoiding} if
$x\notin F(x)$ for every $x\in S$. A subset $H\subseteq S$ is \emph{free} for $F$ if
\[
  x,y\in H,\quad x\ne y\quad\Longrightarrow\quad y\notin F(x).
\]
The question considered here asks whether small pairwise intersections of
the values of a set mapping force a large free set. In the case raised by
Erd\H{o}s and Hajnal \cite[Problem~35]{EH}, one assumes the generalized
continuum hypothesis, takes $\mu=\aleph_\omega$ and $\kappa=\mu^+$, and requires
\[
  f:\kappa\longrightarrow[\kappa]^{\leq\mu},\qquad
  |f(\alpha)\cap f(\beta)|<\mu\quad(\alpha\ne\beta).
\]
We give a counterexample under the weaker cardinal-arithmetic assumption
$2^\mu=\mu^+$. Throughout, we work in ZFC, identify cardinals with their
initial ordinals, and write $[S]^\mu$ for the family of subsets of $S$ of
cardinality $\mu$.

\begin{theorem}\label{thm:main}
Let $\mu$ be an uncountable cardinal with $\cf(\mu)=\aleph_0$ and
$2^\mu=\mu^+$. Put $\kappa=\mu^+$. There is a self-avoiding mapping
$f:\kappa\to[\kappa]^{\leq\mu}$ such that
\[
  |f(\alpha)\cap f(\beta)|<\mu
  \qquad\text{for all distinct }\alpha,\beta<\kappa,
\]
but no subset of $\kappa$ of cardinality $\kappa$ is free for $f$.
\end{theorem}

The intersection bound is pointwise: the construction does not require a
single cardinal below $\mu$ to bound all pairwise intersections. This
distinction is essential at a singular cardinal.

The use of values concentrated near their indexing points is motivated by
Hechler's construction discussed in \cite[p.~279]{EH}. All ingredients
needed for the construction below are proved here.

\section{Proof of the main theorem}

\begin{proof}[Proof of Theorem~\ref{thm:main}]
Fix an increasing sequence $\langle\mu_n:n<\omega\rangle$ of infinite
cardinals below $\mu$ with supremum $\mu$. Put $\kappa=\mu^+$.

\subsection{The sequence space and its enumerations}
Let $X=\mu^\omega$, with the product topology obtained by giving $\mu$ the
discrete topology. For $x\in X$ and $n<\omega$, define
\begin{equation}\label{eq:neighborhoods}
 U_n(x)=\{y\in X:y\mathbin{\upharpoonright}n=x\mathbin{\upharpoonright}n\}.
\end{equation}
These sets form a decreasing neighborhood base at $x$. All closures below
are taken in $X$. If $x\ne y$, then $U_N(x)$ and $U_N(y)$ are disjoint for
some finite $N$. The collection $\mathcal B$ of all finite-prefix cylinders
is a base of cardinality $\mu$, since there are $\mu$ finite sequences from
an infinite alphabet of size $\mu$.

We first verify that $|X|=\kappa$. If $|X|\leq\mu$, list its points as
$\langle z_\xi:\xi<\mu\rangle$, allowing repetitions. For each $n$, choose
\[
 z(n)\in\mu\setminus\{z_\xi(n):\xi<\mu_n\}.
\]
This choice is possible because $\mu_n<\mu$. For every $\xi<\mu$, some $n$
satisfies $\xi<\mu_n$, and then $z(n)\ne z_\xi(n)$. Thus $z$ is absent from
the list, a contradiction. On the other hand,
\begin{equation}\label{eq:cardX}
 |X|=\mu^{\aleph_0}\leq(2^\mu)^{\aleph_0}
      =2^{\mu\cdot\aleph_0}=2^\mu=\kappa.
\end{equation}
Since $\kappa$ is the successor of $\mu$, these inequalities give
$|X|=\kappa$. Fix a bijective enumeration
$X=\{x_\alpha:\alpha<\kappa\}$.

The family $[X]^\mu$ also has cardinality $\kappa$. Indeed,
\begin{equation}\label{eq:cardFamily}
 |[X]^\mu|\leq\kappa^\mu=(2^\mu)^\mu
             =2^{\mu\cdot\mu}=2^\mu=\kappa.
\end{equation}
For the reverse inequality, fix $B\in[X]^\mu$. The sets $B\cup\{x\}$, for
$x\in X\setminus B$, are $\kappa$ distinct members of $[X]^\mu$. We may
therefore enumerate this family as $\langle A_\xi:\xi<\kappa\rangle$.

\subsection{Scheduling the requirements}
For each $\alpha<\kappa$, choose an injection $h_\alpha:\alpha\to\mu$;
this is possible because $|\alpha|\leq\mu$. For $\xi<\alpha$, put
\begin{equation}\label{eq:depth}
 d_\alpha(\xi)=\min\{n<\omega:h_\alpha(\xi)<\mu_n\}.
\end{equation}
For every finite $N$, injectivity of $h_\alpha$ and monotonicity of the
sequence $\langle\mu_n\rangle$ give
\begin{equation}\label{eq:shallow}
 |\{\xi<\alpha:d_\alpha(\xi)<N\}|\leq\mu_N<\mu.
\end{equation}
In fact, $d_\alpha(\xi)<N$ implies $h_\alpha(\xi)<\mu_N$. For $N=0$ the
set on the left is empty.

\subsection{The set mapping}
For each $\alpha<\kappa$, consider the indices
\[
 I_\alpha=\{\xi<\alpha:x_\alpha\in
             \overline{A_\xi\setminus\{x_\alpha\}}\}.
\]
For every $\xi\in I_\alpha$, choose a point
\[
 y_{\alpha,\xi}\in(A_\xi\setminus\{x_\alpha\})
                    \cap U_{d_\alpha(\xi)}(x_\alpha).
\]
The defining closure condition makes this intersection nonempty. Define
\begin{equation}\label{eq:mapping}
 F(x_\alpha)=\{y_{\alpha,\xi}:\xi\in I_\alpha\}.
\end{equation}
Every chosen point differs from $x_\alpha$, so $x_\alpha\notin F(x_\alpha)$.
Moreover,
$|F(x_\alpha)|\leq|I_\alpha|\leq|\alpha|\leq\mu$.
By \eqref{eq:shallow}, for every $\alpha<\kappa$ and every finite $N$,
\begin{equation}\label{eq:exceptions}
 |F(x_\alpha)\setminus U_N(x_\alpha)|\leq\mu_N<\mu.
\end{equation}
Indeed, if $d_\alpha(\xi)\geq N$, then the selected point belongs to
$U_N(x_\alpha)$. Only indices of depth less than $N$ can contribute points
outside that neighborhood. Repeated choices do not increase their number.

\subsection{Pairwise intersections}
Let $\alpha\ne\beta$. Choose a finite $N$ such that
$U_N(x_\alpha)\cap U_N(x_\beta)=\varnothing$. Every common point of the two
values of $F$ lies outside at least one of these neighborhoods. Hence
\[
\begin{split}
 F(x_\alpha)\cap F(x_\beta)\subseteq{}&
   \bigl(F(x_\alpha)\setminus U_N(x_\alpha)\bigr)\\
   &\mathbin{\cup}\bigl(F(x_\beta)\setminus U_N(x_\beta)\bigr).
\end{split}
\]
Applying \eqref{eq:exceptions}, and using that $\mu_N$ is infinite, we obtain
\begin{equation}\label{eq:intersections}
 |F(x_\alpha)\cap F(x_\beta)|
        \leq\mu_N+\mu_N=\mu_N<\mu.
\end{equation}
Only two exceptional sets are combined here. No estimate for a countable
union of sets of size less than the singular cardinal $\mu$ is used. The
depth $N$ may depend on the pair $\alpha,\beta$.

\subsection{No large free set}
Let $H\subseteq X$ have cardinality $\kappa$. For each basic open set
$B\in\mathcal B$ that meets $H$, choose one point of $B\cap H$, and let $D$
be the set of chosen points. Then $D\subseteq H$, $|D|\leq\mu$, and
$H\subseteq\cl(D)$. Indeed, every basic neighborhood of a point of $H$
contains one of the chosen representatives.

Enlarge $D$ inside $H$ to a set $A$ of cardinality exactly $\mu$. Then
$H\subseteq\cl(A)$. Write $A=A_\xi$. Since both $A$ and
$\{x_\alpha:\alpha\leq\xi\}$ have cardinality at most $\mu$, their union
cannot exhaust $H$. Choose
\[
 x_\alpha\in H\setminus A\qquad\text{with}\qquad\alpha>\xi.
\]
Because $x_\alpha\notin A$ and $x_\alpha\in\cl(A)$, we have
\[
 x_\alpha\in\overline{A_\xi\setminus\{x_\alpha\}}.
\]
Thus $\xi\in I_\alpha$, and the construction selected a point
\[
 y_{\alpha,\xi}\in F(x_\alpha)\cap A\subseteq F(x_\alpha)\cap H.
\]
This point is distinct from $x_\alpha$. Hence $H$ is not free for $F$.
Since $H$ was arbitrary, $F$ has no free subset of cardinality $\kappa$.

\subsection{Transfer to the ordinal domain}
Use the fixed bijection $\alpha\mapsto x_\alpha$ to define
\begin{equation}\label{eq:transport}
 f(\alpha)=\{\beta<\kappa:x_\beta\in F(x_\alpha)\}.
\end{equation}
The bijection preserves self-avoidance, the cardinalities of values and
their intersections, and freeness. Therefore $f$ has all the asserted
properties.\qedhere
\end{proof}

\section{The Erd\H{o}s--Hajnal question}

\begin{corollary}\label{cor:GCH}
Assume the generalized continuum hypothesis, and put $\mu=\aleph_\omega$
and $\kappa=\aleph_{\omega+1}$. There is a self-avoiding mapping
$f:\kappa\to[\kappa]^{\leq\mu}$ whose distinct values have intersection of
cardinality less than $\mu$, but which has no free set of cardinality $\kappa$.
\end{corollary}

\begin{proof}
Take $\mu=\aleph_\omega$. This cardinal has cofinality $\aleph_0$, and GCH
gives $2^\mu=\mu^+=\aleph_{\omega+1}$. Theorem~\ref{thm:main} applies.\qedhere
\end{proof}

Identifying $\aleph_{\omega+1}$ with its initial ordinal
$\omega_{\omega+1}$ gives exactly the domain in the Erd\H{o}s--Hajnal
question. Thus the stated pointwise intersection condition does not force
a free set of the required size.

\subsection*{Acknowledgments}
The authors acknowledge the assistance of OpenAI's ChatGPT in developing
the construction and in preparing and checking the mathematical exposition.

\begin{thebibliography}{1}
\bibitem{EH} P.~Erd\H{o}s and A.~Hajnal,
\emph{Unsolved and solved problems in set theory},
in \emph{Proceedings of the Tarski Symposium},
Proceedings of Symposia in Pure Mathematics \textbf{25},
American Mathematical Society, Providence, RI, 1974, pp.~269--287.
\end{thebibliography}
\end{document}
